\chapter{Patterns of Dictatorship}

The defeat and expulsion of the united opposition at the party congress of December 1927 removed the last formidable obstacle in the way of Stalin's progress towards absolute power. Rifts soon appeared in the opposition itself. The attitude of Kamenev at the congress already smacked of surrender. A month later Zinoviev and Kamenev issued a statement that they had ``parted company'' with Trotsky's group, that they rejected its policies, and that their motto now was: ``Back into the party, back into Comintern''. Other defections followed, including some of Trotsky's own followers. The process was hastened when the new turn in official policy became apparent. Trotsky had confidently predicted that the victory of Stalin and Bukharin would presage a sharp reaction to the Right. What happened was exactly the contrary. The grain collections of the first months of 1928 proved that Stalin had abandoned the appeasement of the peasant which the opposition had condemned. Stalin scarcely waited to expel Trotsky from the party and from Moscow before embarking on a policy of forced industrialization at a pace, and at a cost for other sectors of the economy, far beyond anything hitherto contemplated by Trotsky or by anyone else. The exiles languishing in Siberia could now persuade themselves that Stalin had adopted the policies of the opposition, and that their role was to help and support those engaged on the work. This at any rate provided an honourable ground for surrender. Inducements as well as threats were offered to defectors. In June 1928 Zinoviev and Kamenev with some forty other penitents were re-admitted to the party. 
% 注：去除了段首的 "|The", "[apparent." 的单侧括号，以及无意义的乱码 "UvlvlvL"。修复了 "itself The" 中缺少的句号。

Throughout the year Trotsky in Alma-Ata kept up a long-distance correspondence with exiles all over Siberia, and strove, with diminishing success, to strengthen their resistance. He was particularly hurt when Preobrazhensky and Radek, whom he had hitherto counted among his firmest adherents, announced their disagreement with him, and made overtures to the authorities in Moscow. Of the prominent former opposition leaders, only Rakovsky still shared Trotsky's view that Stalin's personal dictatorship and the degeneration of the party were key issues on which no compromise was permissible. Trotsky himself was indefatigable. In the summer of 1928 he sent to the Comintern secretariat a long Critique of the draft programme of Comintern submitted to the congress, which could not be concealed from foreign delegates. It constituted a biting attack on the doctrine of socialism in one country, which was held responsible for all the disasters of Comintern policy. For Stalin, Trotsky, even isolated in a remote corner of the USSR, still represented a focus of dissent, an organized challenge to his authority; and he decided to be rid of him. To hold one of the heroes of the revolution behind bars would at this time have been inconceivable---a measure of the distance which still separated this period from that of the great purges. The problem was to find a destination to which he could be sent. Neither Germany nor any other European country would admit the notorious revolutionary. Turkey, however, proved willing; and in January 1929 Trotsky was conveyed to Odessa and put on board a ship bound for Istanbul. For nearly four years he found refuge on the island of Prinkipo. 
% 注：去除了 "ifor Istanbul" 的 "i"。

Whether or not Stalin under-estimated the damaging effects of Trotsky's indomitable campaign against him in the outside world, it was true that within the USSR he had divested himself of his last serious rival. The groups in the party which thereafter contested his authority offered no threat to his monopoly of power. They did not organize their sympathizers, and, like the united opposition, found it difficult to present any positive alternative programme. Both the united opposition and the later dissidents used traditional language in condemning the evils of bureaucracy and the repression of independent opinion. ``Our disagreements with Stalin'', Bukharin told Kamenev in June 1928, ``are far more serious than those we had with you''. But this was not strictly true. There was an important difference between them, resulting in part from the shift in Stalin's own attitude after his victory at the end of 1927. Trotsky, Zinoviev and Kamenev had criticized Stalin for betraying the aims of the revolution, and compromising with kulaks at home and nationalists and social-democrats abroad; this was an attack from the Left. Bukharin, Rykov and Tomsky blamed the haste and ruthlessness with which Stalin pursued the aims of the revolution, and sought to moderate the pace and the intensity of the pursuit; these were attacks, in the terminology of the period, from the Right. Nor did the later dissidents, like the earlier opposition, automatically exclude themselves from the framework of the party. They were commonly described as guilty, not of opposition, but of ``deviation''. 
% 注：去除了 "|Of Trotsky's", "jhad criticized", "ol the party" (修正为 of) 等边缘字符。

The new group of Right ``deviators'' began to form within a few weeks of the downfall of the united opposition, and well before Trotsky's final banishment from the USSR. Rykov, who had long stood on the party Right, openly expressed his dislike, which was shared by many party members, of the forcible grain collections of January-February 1928. Bukharin was slower to declare himself. He had worked hand-in-glove with Stalin in the campaign against Trotsky. But, once the opposition was defeated, he became expendable, and Stalin soon set to work to undermine his influence. Already at the party congress which expelled the opposition in December, Bukharin incurred sly taunts of neglecting what was called ``the Right danger''. These referred specifically to the affairs of Comintern, but had wider implications. In May 1928, Stalin addressed the Institute of Red Professors, of which Bukharin was the director, deprecating proposals to slow down the progress of industrialization, and speaking of the need to strengthen the Kolkhozy and Sovkhozy, and to improve the grain collections. Though Bukharin's name was not mentioned, the challenge to his views was unmistakable. About the same time Bukharin addressed two memoranda to the Politburo calling in question the pace of industrialization, the pressure which it imposed on the peasant, and the practicability of collective agriculture; and Tomsky became uneasy about the repercussions of industrialization on the workers and on his role in the trade unions. Bukharin, among his other functions, was the editor of the party newspaper, Pravda, and a member of the editorial board of the party journal, Bol'shevik. New appointments were made to the boards of both organs with the evident motive of curbing Bukharin's authority. At the crucial meeting of the party central committee in July, Rykov, Bukharin and Tomsky, all members of the Politburo, appeared as a minority of three to contest current policies. Bukharin, in virtue of his reputation as the leading party theorist and of his skill in debate, emerged as leader of the group. 
% 注：此处并无明显 OCR 错误。

The moment was not yet ripe for an open break. The proceedings ended in an empty compromise, and the show of unanimity in the Politburo was maintained. But Bukharin had now read the signals. While the session was in progress, with the knowledge of Rykov and Tomsky, he secretly visited Kamenev, and offered him a coalition with the remnants of the old opposition against Stalin. Stalin he described as ``a Genghis Khan'', who ``will wait for us to start a discussion and then cut our throats''. It was a belated and futile gesture. The united opposition had been shattered and dispersed; and Kamenev was a broken reed. Bukharin was no tactician. But, when Stalin learned of the demarche some time later, it must have confirmed his determination to crush and humiliate Bukharin. Later in July, Bukharin presided at the sixth congress of Comintern. But Stalin publicly snubbed him by insisting on amendments to the theses which he had drawn up for the congress; and many of the delegates knew that his star was waning. At the end of September he launched his economic broadside ``Notes of an Economist'', and departed on holiday. But he did nothing to organize his supporters, and left the field open for an intensive propaganda campaign against his opinions, though none of the major dissidents were as yet mentioned by name. The session of the party central committee in November once more resulted ostensibly in a compromise (see p. 140 above). But this time it was plainly Bukharin who had executed a retreat to preserve formal unity, and had suffered a crushing defeat. 
% 注：此处并无明显 OCR 错误。

Tomsky proved less pliable, and was the first of the trio to incur public disgrace. He opened the trade union congress of December 1928, a month after the defeat in the party central committee, without attempting to canvass contentious questions. But his reluctance, and that of the other trade union leaders, to face the issue of industrialization, was clear. Pravda accused the trade unions of taking an ``a-political'' line, i.e. of concentrating on the day-to-day interests of the workers, and neglecting ``the new tasks of the reconstruction period''. The Politburo showed its determination to bring Tomsky to heel by appointing Kaganovich, one of Stalin's most militant henchmen, as delegate of the party central committee to the trade union central council. Tomsky indulged in the bold, but empty, gesture---for which he was sharply censured---of resigning his post as president of the trade union central council, and absenting himself from the final session of the congress. Though he retained his membership of the central council for three months longer, he never again appeared on a trade union platform. 
% 注：去除了 "jjline", "'Iworkers", "yperiod", "bf the" 的边缘乱码。

Bukharin did not long survive him. In January 1929, reduced to desperation, he had two more fruitless meetings with Kamenev; Kamenev's record of the earlier meeting was now being circulated in inner party circles. The break could no longer be avoided. It came at a joint session of the Politburo with the party control commission at the end of January. The three dissidents tendered their resignations; and Bukharin made a direct attack, though not by name, on Stalin, protesting against the oppressive regime in the party and ``against the decisions of the party leadership being taken by a single person''. Stalin retorted with a vituperative analysis of the zigzags of Bukharin's record, and of his early disputes with Lenin, and denounced ``the Right opportunist, capitulationist platform'' of the dissidents. The resolution adopted at the end of the session on February 9 rehearsed the catalogue of Bukharin's errors, and convicted him of disloyalty to the party. But it was neither published nor officially communicated to the party central committee, so that Bukharin's status was still formally intact. It was not till April that the central committee met, and, having listened to a further sweeping indictment by Stalin of Bukharin's record, confirmed the resolution of February 9, and removed Bukharin from further work in Pravda and in Comintern, and Tomsky from the trade union central council. But this merely provided formal endorsement of situations which already existed. After the meeting of the central committee, the decisions were communicated by Molotov to the large party conference which had been summoned to approve the first five-year plan. But neither Molotov's report nor the resolution approving the decisions was published. No word of Bukharin's downfall had yet appeared in the press, or percolated to the world at large. This extreme caution was characteristic of Stalin, who did not rate Bukharin as a dangerous opponent, and saw no reason to force this issue. But it was also a tribute to Bukharin's popularity among rank-and-file party members, many of whom, especially in the countryside, shared his moderate inclinations. 
% 注：去除了 "^a direct" 的乱码；修复了断句 ". This extreme caution", "/treason" (修正为 reason) 以及 "['f|popularity", "•^///especially" 等非常杂乱的符号。

The issue next arose when IKKI met in July 1929. At first nobody mentioned Bukharin's absence. But half-way through the session Molotov arrived to deliver a frank denunciation of the three dissidents, and in particular of Bukharin, who had engaged in ``a Right deviation'' and attacked ``our socialist economy''. After this, many delegates, Soviet and foreign, joined in the chorus; and a resolution was passed at the end of the session condemning Bukharin, and approving the decision of the party central committee to debar him from further participation in Comintern and its organs. But once more this resolution was not published with the rest, and appeared in Pravda only some weeks later. Now, however, a thorough-going campaign of denunciation was launched in the press. The climax was reached at a session of the party central committee in November. The three dissidents were induced to sign a somewhat half-hearted recantation of their views, which was published in Pravda. Bukharin was removed from the Politburo, Tomsky and Rykov merely censured and warned not to offend again. By slow process of attrition the dissidents had been discredited, and rendered helpless and harmless. 

A month later, on December 21, 1929, Stalin celebrated his fiftieth birthday. The occasion epitomized tendencies which had grown up insensibly out of his struggles against rivals and his ascent to supreme power. Ever since his appointment, in the last year of Lenin's active life, as general secretary of the party, his strength had lain in his rigid and meticulous management of the party machine, which controlled appointments to key posts in party and state. His approval was a sure avenue to advancement. He had gathered round him a body of faithful henchmen---mostly, party leaders of the second rank---whose political fortunes were linked with his, and who owed him unquestioning personal allegiance. A policy of recruitment inaugurated by the Lenin enrolment of 1924 had built up a rank-and-file membership of reliable workers known for their ready submission to the party line. In the less congenial field of party doctrine Stalin took pains to present himself, not as an innovator, but as the devoted disciple of Lenin and the custodian of party orthodoxy; the spurious attempt to attribute the theory of socialism in one country to Lenin was an example of his eagerness to ground his authority in that of the master. The parallel was consciously cultivated by his entourage. Stalin's words, like those of Lenin, were constantly quoted in the press and in the speeches of his followers, and treated as authoritative. His portrait appeared everywhere in public places, often side by side with that of Lenin. These practices reached their apogee in the birthday tributes, which were marked by a hitherto unprecedented display of personal adulation. 
% 注：去除了 ",a sure avenue" 中的逗号。

Many traits, however, distinguished the character of Stalin's rule from anything that could have been imagined under Lenin. Stalin had a form of vanity totally alien to Lenin, which demanded, not indeed the holding or the trappings of office, but absolute obedience and the recognition of his infallibility. No overt criticism, no expression of dissent, appeared any longer in the party press, or even in specialist journals. Such discussions of current questions as could still be found were marked by tasteless and uniform tributes to the leader, and by the celebration of often mythical achievements. Stalin became a remote, isolated figure, exalted far above ordinary mortals, and indeed above his closest colleagues. He seems to have lacked any warmth of feeling for his fellow-men; he was cruel and vindictive to those who thwarted his will, or excited his resentment or his antipathy. His commitment to Marxism and socialism was only skin-deep. Socialism was not something that grew out of the objective economic situation and out of the revolt of class-conscious workers against the oppressive domination of capitalism; it was something to be imposed from above, arbitrarily and by force. Stalin's attitude to the masses was contemptuous; he was indifferent to liberty and to equality; he was scornful of the prospects of revolution in any country outside the USSR. He was the only member of the party central committee who as early as January 1918 had maintained, in opposition to Lenin, that ``there is no revolutionary movement in the West''. The commitment to socialism in one country, though the attitudes which crystallized into the new doctrine were not exclusively of Stalin's making, perfectly fitted the man. It enabled him to match professions of socialism with Russian nationalism, the only political creed which moved him at all deeply. In Stalin's treatment of national minorities or of the smaller nations, nationalism easily degenerated into chauvinism. Notes of the old Russian anti-Semitism, sternly denounced by Lenin and the early Bolsheviks, were heard; and official condemnations of it, though persistent, began to sound less decisive. In art and literature, the eager experimentalism of the first years of the revolution was abandoned in favour of a return to traditional Russian models, enforced by an increasingly strict censorship. Marxist schools of history and law passed under a cloud; to seek continuity with the Russian past was no longer a cause for reproach. Socialism in one country pointed back to an old Russian national exclusiveness, rejected by Marx as well as by Lenin. It was not altogether incongruous to place Stalin's regime in the context of Russian history. 
% 注：去除了 "jlappeared", "//journals", "f/imen", "|his commit- tment", "iiwas", "/situation", "/the oppressive", "'to be imposed", "I attitude" 等等极其密集的左侧竖向噪点。

This constriction of the revolution within a narrow nationalist strait-jacket had its reverse side. It would be unfair to depict Stalin as a man moved exclusively by the lust for personal power. He dedicated his indefatigable energy to the transformation of primitive peasant Russia into a modern industrial Power, capable of facing the major capitalist Powers on equal terms. The need to ``catch up with'' or ``overtake'' the capitalist countries was an obsessive theme, and inspired most of the rare purple passages in Stalin's drab prose. It provided the peroration of his anniversary article of November 1929 on ``The Year of the Great Break-Through'': 

\begin{quote}
We are going full steam along the road of industrialization to socialism, leaving behind our age-long ``Russian'' backwardness. We are becoming a country of metal, a country of the automobile, a country of the tractor. And when we have seated the USSR in an automobile, and the peasant on a tractor, then let the honourable capitalists, who plume themselves on their ``civilization'', try to catch us up. We shall see which countries can then be counted as backward and which as advanced. 
\end{quote}
% 注：将这三段被大量符号干扰的引文合并为一个 quote。去除了 "jjl We", "///to", "/ It wardness", "llP>", "//hutomobile", "y// And", "/rfhnd", "/ /Icapitalists", "(/try" 等极端的乱码。

And later, rather more soberly, he conjured up a picture of Russia through the ages, ``beaten for her backwardness'' by a succession of foreign invaders, from ``Mongol khans'' to ``Anglo-French capitalists'' and ``Japanese barons'', and concluded: 

\begin{quote}
We have lagged 50 or 100 years behind the advanced countries. We must close this gap in ten years. Either we shall do it, or they will crush us. 
\end{quote}
% 注：同上，去除了 "jjl We", "///countries", "f'lwe" 等乱码。

This extraordinary combination of a commitment to the industrialization and modernization of the economy, which appealed to convinced Marxists as a vital step on the road to socialism, and a commitment to a revival of the power and prestige of the Russian nation, which appealed to the army, to the bureaucratic and technological elites, to all survivors of the old regime who had entered the service of the new, gave Stalin his unbreakable hold over the party, the government and the administration. To attribute this simply to Stalin's political astuteness, or to the efficiency of the machine, or to the severity of measures taken to suppress dissent, would be a mistake. It was not only defectors from the opposition in 1928 and 1929 who felt that Stalin's unflinching determination in the pursuit of long-desired ends outweighed the harsh methods used to enforce his policies. Some people argued that without these methods the ends could not be attained, others that they could not be attained without Stalin's forceful leadership, and that it was therefore necessary to tolerate his unwelcome idiosyncracies. The fact that this was a revolution from above, and that it placed the heaviest burdens on the very classes which were its declared beneficiaries, did not seriously disturb the picture. Enthusiasm for the great leap forward engulfed many party members and others engaged, in one capacity or another, in forwarding the grand design, and left them indifferent to other considerations. It was a society well accustomed to associate government with oppression, and to treat it as an inescapable evil. 
% 注：去除了 "^the grand", "govern- merit" (修正为 government) 等乱码和断字。

Stalin on his fiftieth birthday stood at the summit of his ambition. Enough had indeed occurred to lend weight to Lenin's apprehensions of his rough and arbitrary use of power. He had already displayed an extraordinary ruthlessness in enforcing his will and in crushing all opposition to it. But the full revelation of the quality of his dictatorship still lay ahead. The horrors of the process of collectivization, of the concentration camps, of the great show trials, and of the indiscriminate killing, with or without trial, not only of those who had opposed him in the past, but of many who had assisted his rise to power, accompanied by the imposition of a rigid and uniform orthodoxy on the press, on art and literature, on history and on science, and by the suppression of every critical opinion, left a blot which could not be erased by victory in the war or by the sequel. The fluctuations in Stalin's reputation since his death in the eyes of his compatriots seem to reflect confused and conflicting emotions of admiration and shame. This ambivalence may persist for a long time. The precedent of Peter the Great has often been invoked, and is astonishingly apposite. Peter too was a man of formidable energy and extreme ferocity. He revived and outdid the worst brutalities of earlier Tsars; and his record excited revulsion in later generations of historians. Yet his achievement in borrowing from the west, in forcing on primitive Russia the material foundations of modern civilization, and in giving Russia a place among the European Powers, obliged them to concede, however reluctantly, his title to greatness. Stalin was the most ruthless despot Russia had known since Peter, and also a great westernizer. 
% 注：去除了 "I lambition", "//Lenin's", "'ll He", "/(/enforcing", ":", "/triots", "i i", "_.)1" 等大量残破的边缘标记。